LLVManim bridges the gap between two observations:
\begin{enumerate}
    \item Methods exist to quickly generate code visualizations, but they produce
    output that is unappealing for presentation purposes.
    \item The Manim library can produce visually impressive animations, but using
    it requires significant manual effort.
\end{enumerate}
By automating the translation from LLVM IR to Manim scenes, educators
and tech leads alike can produce presentation-ready visualizations
without needing to compromise between style and accessibility.

\needspace{6\baselineskip}

\subsection{Design Goals}
\label{sec:llvmanim:design_goals}
\subimport{subsec/}{design_goals.tex}

\begin{figure*}[htbp]
\includegraphics[width=\textwidth]{05_cli_dispatch}
\caption{CLI dispatch flow. Each command flag (\texttt{-{}-json}, \texttt{-{}-draw}, \texttt{-{}-animate}, \texttt{-{}-cfg-animate}) shares the Ingest pipeline and branches to the appropriate Transform and Render stages.}
\label{fig:cli_dispatch}
\end{figure*}

\needspace{8\baselineskip}
\subsection{Architecture}
\label{sec:llvmanim:architecture}
\subimport{subsec/}{architecture.tex}
